\documentclass[12pt]{article}

\addtolength{\textwidth}{.8in}
\addtolength{\hoffset}{-.4in}
\addtolength{\textheight}{1in}
\addtolength{\voffset}{-1in}

\usepackage{amsmath}
\usepackage{graphicx}
\usepackage{upquote} % to convert funny quotes to straight quotes
\usepackage{amsfonts}
\usepackage{paralist}
\setlength{\parindent}{15pt}

\renewcommand*{\thefootnote}{\fnsymbol{footnote}}


\begin{document}

\begin{center}
\begin{Large}
DATA 219 Notes -- Feb. 20, 2018
\end{Large}
\end{center}

\section{Time is complicated}

One of the most common categories of information you'll ever deal with are
those dealing with various aspects \textit{time}. It turns out this is
considerably more complex than you might think at first. For one thing, the
English word ``time" can mean several different things -- a point in time
(3:05pm on Feb.~20, 2018), an abstract wall clock time (3:05pm), a duration of
time (2 and a half hours), \textit{etc.} For another, the clock and calendar
systems that humans have devised are just plain weird. (Why does it go from
11am to 12\underline{\textbf{p}}m to 1pm? Why do some months have more days
than others? Why do some years have more than others? Why doesn't every month
start over on a Monday? \textit{Etc.} There are good answers to some of these
questions, and lousy answers to others, but they all make dealing with time a
major pain.

Python helps deal with the pain through the \texttt{datetime} package. Here's
the stuff we need:

\begin{itemize}
\itemsep1em

\item \textbf{The \texttt{timedelta} type}. Objects of this type represent
\textit{durations in time}. You could create one like this:
    \begin{itemize}
    \itemsep.1em
    \item \texttt{one\_semester = datetime.timedelta(weeks=15)}
    \item \texttt{this\_class = datetime.timedelta(hours=1, minutes=15)}
    \item \texttt{finger\_snap = datetime.timedelta(seconds=.1)}
    \end{itemize}

    (You can also use \texttt{days}, \texttt{milliseconds}, and
\texttt{microseconds}. Anything else is ambiguous.)

    Don't think of the \texttt{timedelta} object you create as being
intrinsically ``in" any particular unit -- just think of it as an ``amount of
time," which could be expressed in any units.

    \texttt{timedelta}s can be positive or negative, and you can add or
subtract them to/from each other.

\item \textbf{The \texttt{date} type}. Each \texttt{date} object represents a
date in an idealized Gregorian calendar. You can create one via:
    \begin{itemize}
    \itemsep.1em
    \item \texttt{sunny\_day = datetime.date(2018,2,20)}
    \item \texttt{here\_we\_are = datetime.date.today()}
    \end{itemize}

    You can \textit{subtract} two \texttt{date}s to get a \texttt{timedelta}.
You can \textit{add} a \texttt{date} to a \texttt{timedelta} to get another
\texttt{date}. Note that if your \texttt{timedelta} has precision more
granular than days, it is lost in these operations:

\small
\begin{Verbatim}[fontsize=\small,samepage=true,frame=none]
>>> datetime.date.today() == datetime.date.today() + finger
True

>>> datetime.date.today() == datetime.date.today() + this_class
True

>>> datetime.date.today() == datetime.date.today() + one_semester
False
\end{Verbatim}
\normalsize

    You can also compare dates with comparison operators like \texttt{<} and
\texttt{>=}.

\item \textbf{The \texttt{time} type.} A ``time," by itself, represents an
abstract wall clock time, independent of any particular day. It doesn't mean
``9am this morning" but rather ``9am's in general."
    \begin{itemize}
    \itemsep.1em
    \item \texttt{start\_of\_class = datetime.time(hours=14, minute=0)}
    \end{itemize}

    These can also be compared in expected ways, but \textit{not} added or
subtracted.

\item Finally, we have \textbf{the \texttt{datetime} type}, which as its name
suggests, combines features of both dates and times. In fact, a
\texttt{datetime} object represents \textit{an absolute, specific point in
history}. Creation options include:
    \small
    \begin{itemize}
    \itemsep.1em
    \item \texttt{this\_exact\_moment = datetime.datetime.now()}
    \item \texttt{terrorist\_attack = datetime.datetime(2001, 9, 11, 10, 33)}
    \item \texttt{todays\_class = datetime.datetime.combine(here\_we\_are, start\_of\_class)}
    \end{itemize}
    \normalsize

    \texttt{datetime}s can be subtracted from each other, and
\texttt{timedelta}s can be added or subtracted to/from them to produce other
\texttt{datetime}s. They can also be compared. Finally, \texttt{.time()} and
\texttt{.date()} methods will produce the constituent parts as \texttt{time}
or \texttt{date} objects, accordingly.

\item Oh, and one more thing...I should add that there's such a thing as a
``\texttt{timestamp}", which in Python isn't a class, but simply a word
describing a \texttt{float} value. A timestamp is -- get this -- \textit{the
number of seconds that have elapsed since the stroke of midnight Jan.~1, 1970
UTC} (Coordinated Universal Time -- the Greenwich Mean
Timezone).\footnote{That particular instant in time is sometimes called ``The
Epoch." It's an arbitrary reference point, but was chosen in part to be
approximately equal to Stephen's original birthday.} This will occasionally
come up, since even though it's gibberish to humans, it is extremely easy to
work with mathematically.

\end{itemize}

\vspace{.2in}

\subsection{Parsing and formatting}

Now most of the time we won't be creating \texttt{datetime} objects using the
techniques above. Instead, we'll be \textit{parsing} them from string data.

There are a seemingly unlimited number of ways that dates and times can be
expressed in strings. Just consider a few possibilities:

\begin{compactitem}
\item Feb. 20, 2018 2pm
\item Feb 20, 2018 2:00pm
\item Feb 20, 2018 2:00 P.M.
\item 2/20/18 2:00 P.M.
\item 2/20/18 14:00
\item 02/20/18 14:00:00
\item 20-Feb-2018 2:00pm
\item 20-Feb-2018 14:00
\item 2018-Feb-20 14:00:00
\item ...
\end{compactitem}

and on it goes. It's a testimony to the power of the human mind that we can
look at these and at a glance infer what date/time was intended. Getting a
computer to automatically do so is somewhat of a challenge.

\subsubsection{Parsing with \texttt{.strptime()}}

The \texttt{date}, \texttt{time}, and \texttt{datetime} classes all have a
(somewhat clunky) method called ``\texttt{.strptime()}" that you can call,
passing (1) a string representing a date, time, or datetime, and (2) a
formatting string telling the method how to interpret it. This is best seen
through an example:

\begin{Verbatim}[fontsize=\small,samepage=true,frame=none]
my_string = "2/20/2018 14:00"
start_of_class = datetime.datetime.strptime(my_string, "%m/%d/%Y %H:%M")

>>> print(start_of_class)
2018-02-20 14:00:00
\end{Verbatim}

This example demonstrates a full ``round-trip" operation: we started with a
\texttt{string}, parsed it into a \texttt{datetime}, and then printed it back
to the screen as a \texttt{string}. Note that the format of the printout
didn't match the format of ``\texttt{my\_string}"; this is okay. What matters
internally to our program is that our date \& time objects \textit{represent}
the right things, not that they are printed in certain ways.

Now what are all those funky percent signs and junk in the second argument to
\texttt{.strptime()}? They all have a special meaning. Here are some common
ones (others are in the docs):

\begin{center}
\begin{tabular}{|c|l|}
\hline
\%a & Weekday (abbreviated): Sun, Mon, \textit{etc.} \\
\%A & Weekday (full): Sunday, Monday, \textit{etc.} \\
\%w & Weekday (number): 0, 1, ..., 6 \\
\%d & Day of month (number): 0, 1, ..., 31 \\
\%b & Month (abbreviated): Jan, Feb, \textit{etc.} \\
\%B & Month (full): January, February, \textit{etc.} \\
\%m & Month (number): 01, 02, ..., 12 \\
\%y & \textbf{Two}-digit year: 00, 01, ..., 99 \\
\%Y & \textbf{Four}-digit year: ..., 1999, 2000, 2001, ...\\
\hline
\%H & Hour on \textbf{24-hour} clock: 01, 02, ..., 23 \\
\%I & Hour on \textbf{12-hour} clock: 01, 02, ..., 12 \\
\%p & AM/PM (only used in conjunction with \textbf{\%I}) \\
\%M & Minute: 01, 02, ..., 59 \\
\%S & Second: 01, 02, ..., 59 \\
\hline
\end{tabular}
\end{center}

You have to be really careful here and just go slow. Examine the date and/or
time input that you have in your data set, and work out which bits of info,
combined with which literal characters (like ``\texttt{/}" and ``\texttt{:}"
in the example above), appear in which order.

\subsubsection{Formatting with \texttt{.strftime()}}

Complementarily to \texttt{.strptime()} is ``\texttt{.strftime()}". (Note that
the ``p" stands for ``\textbf{p}arse," while
the ``f" stands for ``\textbf{f}ormat.") This is the reverse operation: given
a \texttt{date}, \texttt{time}, or \texttt{datetime} \textit{object}, spit it
out as a string of characters. This isn't used nearly as often, but it
occasionally is, and it's useful to know as a translation mechanism.

\begin{Verbatim}[fontsize=\small,samepage=true,frame=none]
>>> print(start_of_class.strftime("%A, %B the %dth at %I%p."))
Tuesday, February the 20th at 02PM.
\end{Verbatim}

One common error: note that you call \texttt{.strptime()} directly on the
\textit{class} (\texttt{date}, \texttt{time}, or \texttt{datetime}), passing
two arguments:

\begin{Verbatim}[fontsize=\small,samepage=true,frame=none]
an_object = datetime.datetime.strptime(my_birthday_string, "%m/%d/%Y %H:%M")
\end{Verbatim}

but you call \texttt{.strftime()} on the object in question
(a \texttt{date}, \texttt{time}, or \texttt{datetime}), passing
one argument:

\begin{Verbatim}[fontsize=\small,samepage=true,frame=none]
a_string = my_birthday_object.strftime("%m/%d/%Y %H:%M")
\end{Verbatim}


